\documentclass[12pt,oneside]{book}
\usepackage{times,mathptmx}
\usepackage[pdftex]{graphicx}
\usepackage{calc}
\usepackage{tabularx,ragged2e,booktabs,caption,subcaption}
\usepackage{array}
\newcolumntype{L}[1]{>{\raggedright\let\newline\\\arraybackslash\hspace{0pt}}m{#1}}
\newcolumntype{C}[1]{>{\centering\let\newline\\\arraybackslash\hspace{0pt}}m{#1}}
\newcolumntype{R}[1]{>{\raggedleft\let\newline\\\arraybackslash\hspace{0pt}}m{#1}}
\usepackage{multirow}
\usepackage{tocloft}
\usepackage{xcolor}
\usepackage{color,soul}
\usepackage{amsmath}
\definecolor{linknavy}{rgb}{0,0,0.50196}
\definecolor{linkred}{rgb}{1,0,0}
\definecolor{linkblue}{rgb}{0,0,1}
\usepackage{float}
\usepackage{graphpap}
\usepackage{rotating}
\usepackage{graphicx}
\usepackage{geometry}
\usepackage{relsize}
\usepackage{ltablex}
\usepackage{longtable}
\usepackage{lscape}
\usepackage{amssymb}
\usepackage{makeidx} % Create index at end of document
\usepackage[nottoc,notlof,notlot]{tocbibind} % Put the bibliography and index in the ToC
\usepackage{lastpage} % Automatic last page number reference.
\usepackage[T1]{fontenc}
\usepackage{enumerate}
\usepackage{upquote}
\usepackage{moreverb}
\usepackage{xfrac}
\usepackage{cite}
\usepackage{tikz}
\usepackage{fancyhdr}
\usepackage{placeins}
% \usepackage{subfig}
% \usepackage{caption}
\usepackage[toc,page]{appendix}

\usepackage{notoccite}
\newcommand\notype[1]{\unskip}

\usepackage{titlesec}
\titleformat{\chapter}[hang] 
{\normalfont\huge\bfseries}{\chaptertitlename\ \thechapter}{1em}{} 
\titlespacing*{\chapter}{0pt}{-30pt}{20pt}

\newcommand{\nopart}{\expandafter\def\csname Parent-1\endcsname{}} % To fix table of contents in pdf.

\usepackage{siunitx}
\sisetup{
    detect-all = true,
    input-decimal-markers = {.},
    input-ignore = {,},
    inter-unit-product = \ensuremath{{}\cdot{}},
    multi-part-units = repeat,
    number-unit-product = \text{~},
    per-mode = fraction,
    separate-uncertainty = true,
}

\usepackage{listings}
\usepackage{textcomp}
\definecolor{lbcolor}{rgb}{0.96,0.96,0.96}

\usepackage[pdftex,
        colorlinks=true,
        urlcolor=linkblue,     % \href{...}{...} external (URL)
        citecolor=linkred,     % citation number colors
        linkcolor=linknavy,    % \ref{...} and \pageref{...}
        pdfproducer={pdflatex},
        pdfpagemode=UseNone,
        bookmarksopen=true,
        plainpages=false,
        verbose]{hyperref}

\setlength{\textwidth}{6.5in}
\setlength{\textheight}{9.0in}
\setlength{\topmargin}{0.in}
\setlength{\headheight}{0.pt}
\setlength{\headsep}{0.in}
\setlength{\parindent}{0.0in}
\setlength{\itemindent}{0.25in}
\setlength{\oddsidemargin}{0.0in}
\setlength{\evensidemargin}{0.0in}
% \setlength{\leftmargini}{\parindent} % Controls the indenting of the "bullets" in a list
\setlength{\cftsecnumwidth}{0.38in}
\setlength{\cftsubsecnumwidth}{0.5in}
\setlength{\cftfignumwidth}{0.45in}
\setlength{\cfttabnumwidth}{0.45in}
\setlength{\parskip}{1em}

\newcommand{\titlesigs}
{
\medskip
\flushright{Underwriters Laboratories Inc.}\\
{\em Terrence Brady, President}
\medskip
\flushright{UL Firefighter Safety Research Institute\\
{\em Stephen Kerber, Director} \\
\hspace{1in} \\
}
}

\newcommand{\headerB}[1]{
\flushleft{
\fontsize{28}{33.6}\selectfont
\bf{#1}
}
}

\newcommand{\headerC}[1]{
\vspace{.5in}
\flushright{\fontsize{14}{16.8}\selectfont
#1}
}

% \newcolumntype{L}{>{\centering\arraybackslash}m{4cm}}

\floatstyle{boxed}
\newfloat{notebox}{H}{lon}
\newfloat{warning}{H}{low}

\newenvironment{conditions}
  {\par\vspace{\abovedisplayskip}\noindent\begin{tabular}{>{$}l<{$} @{${}={}$} l}}
  {\end{tabular}\par\vspace{\belowdisplayskip}}


% Rename chapter headings
\renewcommand{\chaptername}{}
\renewcommand{\bibname}{References}

\pagestyle{fancy}
\lhead{}
\rhead{}
\chead{}
\renewcommand{\headrulewidth}{0pt}

% COMMENT TO REMOVE WATERMARK
\usepackage{draftwatermark}
\SetWatermarkText{DRAFT}
\SetWatermarkScale{1}

\begin{document}
\pagenumbering{gobble}
\bibliographystyle{unsrt}
	
\begin{minipage}[t][9in][s]{6.25in}

\headerB{
Post-Fire Exposures: Examination of Hot, Warm and Cold Scenes  \\
}

\normalsize
\headerC{
	\flushleft{
	Gavin Horn \\
	Daniel Madrzykowski \\

\vspace{0.2in}
	UL Firefighter Safety Research Institute \\
	Columbia, MD 20145 \\
	\vspace*{2\baselineskip}
	

	Summer Neumann     \\

\vspace{0.2in}
	UL Environment \& Sustainability \\
	Lake Forest, CA 92630  \\
	\vspace*{2\baselineskip}
	}
}
	\vfill

\begin{flushright}
\begin{minipage}{0.5\textwidth}
\includegraphics[width=\textwidth]{UnderwritersLaboratories_TM} \\ 
\end{minipage} \hfill
\begin{minipage}{0.25\textwidth}
\includegraphics[width=2.0in]{FSRI_GraphicShield} \\ 
\end{minipage}
\end{flushright}

\scriptsize
\flushright{Copyright 2020. \textcopyright Underwriters Laboratories Inc. All rights reserved. UL and the UL logo are trademarks of UL LLC}
\end{minipage}

\newpage
\hspace{5in}
\newpage

\begin{minipage}[t][9in][s]{6.25in}
\pagenumbering{gobble}

\headerB{
Post-Fire Exposures: Examination of Hot, Warm and Cold Scenes \\
}

\headerC{
	\flushleft{
	Gavin Horn \\
	Daniel Madrzykowski \\

\vspace{0.2in}
	UL Firefighter Safety Research Institute \\
	Columbia, MD 20145 \\
	\vspace*{2\baselineskip}
	

	Summer Neumann     \\

\vspace{0.2in}
	UL Environment \& Sustainability \\
	Lake Forest, CA 92630  \\
	\vspace*{2\baselineskip}
	}
}
	\vfill

\begin{flushright}
\begin{minipage}{0.5\textwidth}
\includegraphics[width=\textwidth]{UnderwritersLaboratories_TM} \\ 
\end{minipage} \hfill
\begin{minipage}{0.25\textwidth}
\includegraphics[width=2.0in]{FSRI_GraphicShield} \\ 
\end{minipage}
\end{flushright}
\titlesigs

\scriptsize
\flushright{Copyright 2020. \textcopyright Underwriters Laboratories Inc. All rights reserved. UL and the UL logo are trademarks of UL LLC}
\end{minipage}

\newpage

\begin{minipage}[t][9in][s]{6.25in}

\flushleft{In no event shall UL be responsible to anyone for whatever use or non-use is made of the information contained in this Report and in no event shall UL, its employees, or its agents incur any obligation or liability for damages including, but not limited to, consequential damage arising out of or in connection  with the use or inability to use the information contained in this Report. Information conveyed by this Report applies only to the specimens actually involved in these tests. UL has not established a factory Follow-Up Service Program to determine the conformance of subsequently produced material, nor has any provision been made to apply any registered mark of UL to such material. The issuance of this Report in no way implies Listing, Classification or Recognition by UL and does not authorize the use of UL Listing, Classification or Recognition Marks or other reference to UL on or in connection with the product or system.
}

\vspace{3in}
\vfill
\hspace{1in}

\end{minipage}

\newpage

\frontmatter

\pagestyle{plain}
\pagenumbering{roman}

\cleardoublepage
\phantomsection
\tableofcontents

\cleardoublepage
\phantomsection
\addcontentsline{toc}{chapter}{List of Figures}
\listoffigures

\cleardoublepage
\phantomsection
\addcontentsline{toc}{chapter}{List of Tables}
\listoftables

\chapter{List of Abbreviations}

\begin{tabbing}
\hspace{1.5in} \= \\
UL 	    \> Underwriters Laboratories \\
UL FSRI \> UL Firefighter Safety Research Institute \\
\end{tabbing}

\mainmatter

\chapter*{Acknowledgments}
\pagenumbering{gobble}

\begin{table}[!ht]
	\centering
	\caption*{Generic Table}
	\begin{tabular}{ll}
		\toprule[1.5pt]
		Name & Awesomeness \\ 
		\midrule
		Neil Armstrong		 & First person to walk on the moon \\
		Edith Clarke		 & First female ECE professor \\
		John Goodenough		 & Rechargeable Lithium-Ion battery \\
		Nikola Telsa		 & Modern alternating current \\  
		\bottomrule[1.25pt]
	\end{tabular}
\end{table}

\newpage

\chapter*{Abstract}

\newpage
\chapter{Introduction}
\label{chap:intro}
\pagenumbering{arabic}
\setcounter{page}{1}

Fire investigators arrive on the fire scene after the fire is extinguished and the IDLH hazards are mitigated.  However, there are still sufficient physical and chemical hazards at the fire scene. 
 <<Fire  examining a scene several days after the fire, when the air seems safe.>>

The IAAI has defined catagories of fire scenes based on the post-fire extinguishment time and has recommended PPE protection for each time period based on the potential hazard that may exist~\cite{IAAI:2020}.   

\begin{enumerate}
\item HOT SCENE A – A fire scene where the fire has been extinguished, but overhaul has not yet commenced or is in progress.

\item HOT SCENE B – A fire scene that has been fully extinguished/overhauled for less than two hours.

\item  WARM SCENE – A fire scene that has been fully extinguished at least two hours but for less than 72 hours. 

\item COLD SCENE – A fire scene that has been fully extinguished for at least 72 hours and not generating detectable or visible dust, fumes, mists, particulates, gases, vapors or aerosols.

\end{enumerate}




\section{Objectives}

To improve the health and safety of fire investigators through understanding the risk for solid phase and vapor phase chemical exposures at the post-fire scene.  These experiments will characterize air-borne contaminants that may be encountered while investigating a residential fire scene and how these contaminants may evolve in the days following the fire, from within an hour after fire extinguishment to 5 days after the fire.          


\chapter{Measurement Approach}

Area monitoring devices are set up either in the kitchen/living room area or a bedroom depending on the area of origin of the fire.  Samples for polycyclic aromatic hydrocarbons (PAHs), aldehydes, carbon monoxide, hydrogen sulfide, hydrogen cyanide, volatile organic compounds (VOCs -particularly benzene, toluene, ethylbenzene and xylene (BTEX)), fibers and a variety of metals are being collected.  

\section{Instrumentation}
\label{sec:instrument}

\subsection{Measurement Locations}
\label{subsec:measure_locs}

\subsection{Measurement Uncertainty}
\label{subsec:uncertainty}

\textit{\hl{Note: this template contains uncertainty information for a variety of different measurement types. Be sure to remove any types that aren't mentioned.}} 

\textit{\hl{Additionally, although this section provides information about measurement uncertainty in paragraphs that are sufficient for a technical report,}\textbf{\hl{ it's highly recommended that the content obtained from this section is still reviewed after being copied to a technical report to verify the accuracy and applicability of the information.}}\hl{For example, it should be verified that the instrumentation described in the paragraphs below resemble the same type used during the experiments mentioned in the technical report.}}

There are different components of uncertainty in the measured values reported in this document, specifically gas temperature, total heat flux, pressure, gas velocity, gas concentration, heat release rate, length, mass, structure leakage, and water flow rate. Uncertainties are grouped into two categories according to the method used to estimate them. Type A uncertainties are those evaluated by statistical methods, and Type B are those evaluated by other means~\cite{Taylor&Kuyatt:1994}. Type B analysis of systematic uncertainties involves estimating the upper ($+$~a) and lower ($-$~a) limits for the quantity in question such that the probability that the value would be in the interval ($\pm$~a) is essentially 100~\%. After estimating uncertainties by either Type A or B analysis, the uncertainties can be combined in quadrature to yield the combined standard uncertainty. Multiplying this combined standard uncertainty by a coverage factor of two results in an expanded uncertainty with a 95~\% confidence interval (2$\sigma$). For some components, such as the zero and calibration elements, uncertainties were derived from referenced instrument specifications. For other components, referenced research results and past experience with the instruments provided input for the uncertainty determination.

\subsubsection*{Length}
Length measurements, such as the room dimensions and instrumentation locations, were made with either a hand held laser measurement device having an accuracy of 0.25~in. ($\pm$~6.0~mm) over a range of 2~ft (0.6~m) to 50~ft (15.2~m)~\cite{StanleyTools} or $\pm$~0.02~in. ($\pm$~0.51~mm) resolution steel measuring tapes manufactured in compliance with NIST Manual 44~\cite{Butcher:2012}, which specifies a tolerance of $\pm$~0.06~in. ($\pm$~1.5~mm) for 30~ft (9.1~m) tapes and $\pm$~0.25~in. ($\pm$~6.4~mm) for 100~ft (30.5~m) tapes. These uncertainties are all well within the precision of the reported dimensions, which are typically rounded to the nearest inch. Some issues, such as levelness of the device and ``soft'' edges on upholstered furniture, result in an estimated expanded uncertainty of $\pm$~1.0~\% for reported length measurements.

\subsubsection*{Mass}
The load cell used to weigh the fuels prior to the experiments had a range of 0~lb (0~kg) to 441~lb (200~kg) with a resolution of 0.11~lb (0.05~kg) and a calibration uncertainty within 1~\%~\cite{Ohaus:2000}. The total expanded uncertainty for the fuel weights measured by the load cell that are presented in this report is estimated to be less than $\pm$~5~\%.


\chapter{Experimental Configuration}
\label{chap:exp_config}

\section{Experimental Structure}





<<<<<<< HEAD
\section{Structure Contents}


\subsection{Open Area}



\subsection{Bedrooms}




=======
\subsubsection*{Structure Leakage} 
To characterize ventilation within the structure, an air leakage measurement system (Model~5101) was used to measure the amount of leakage associated with the training prop before each test~\cite{retrotec:leakage}. \textit{ASTM~E779-10, Standard Test Method for Determining Air Leakage Rate by Fan Pressurization} was followed to determine the air leakage rate of the prop before each experiment~\cite{astm_e779}. The measured leakage rates were recorded in units of air changes per hour at 50~Pa (ACPH50). Retrotec, the manufacturer of the leakage measurement system, reports an accuracy of $\pm$~5~\% for the system.

\subsubsection*{Water Flow Rate}
Water flow rate was measured with a 1.5~in. (3.8~cm) diameter electromagnetic flow meter from Badger Meter, Inc. (Model~M2000). The meter consisted of stainless steel pipe lined with a non-conductive material. Energized coils on the outside of the non-conductive material imposed a magnetic field across the pipe, and when the conductive fluid (water) flowed across the magnetic field, a voltage proportional to flow velocity was produced. The manufacturer reports a $\pm$~0.25~\% calibration uncertainty for the accuracy of the measurement~\cite{Badger:2015}.

\subsubsection*{Airborne Contaminants}
\hl{Will depend upon technique used to measure each class of chemicals and particulate dimensions.} 

The estimated minimum detectable concentration of PAHs, HCN, and benzene  is  0.0021 mg/m3.

The estimated minimum detectable concentration of acid gas concentrations is 0.175 mg/m3.

\begin{table}[!ht]
	\centering
	\caption{Airborne contaminant experimental measurements}
	\begin{tabular}{lccc}
		\toprule[1.5pt]
		Compound & Method & Media & Pump \\
		\midrule
		<row\_1\_content>     & <row\_1\_content>     & <row\_1\_content>     & <row\_1\_content>     \\
		<row\_2\_content>     & <row\_2\_content>     & <row\_2\_content>     & <row\_2\_content>     \\
		\bottomrule[1.25pt]
	\end{tabular}
\end{table}


\section{Fuel Load}

\chapter{Experimental Procedure}
\label{chap:exp_procedure}

The following structure should be used to define tables:
\begin{verbatim}
	\begin{table}[!ht]
	    \centering
	    \caption{Example of Style in which Table Caption Should be 
        Written}
	    \label{tab:<ref_name>}
	    \begin{tabular}{lc}
	        \toprule[1.5pt]
	        <header\_1> & <header\_2> \\
	        \midrule
	        <row\_1\_content>     & <row\_1\_content>     \\
	        <row\_2\_content>     & <row\_2\_content>     \\
	        \bottomrule[1.25pt]
	    \end{tabular}
	\end{table}
\end{verbatim}


\begin{table}[!ht]
	\centering
	\caption{Example of Style in which Table Caption Should be Written}
	\label{tab:<ref_name>}
	\begin{tabular}{llll}
		\toprule[1.5pt]
		Item & Dimensions (in) & Materials & Mass (kg)  \\
		\midrule
		Mattress Topper  & 75 x 58 x 4  & Polyurethane Foam  &  6.85     \\
		Mattress     & 79 X 59 x12  & 90~\% Polyurethane foam, 10~\% Blended rayon \& polyester   & 30.1  \\
		<row\_1\_content>     & <row\_2\_content>  & <row\_3\_content>     & <row\_4\_content>     \\
		<row\_1\_content>     & <row\_2\_content>  & <row\_3\_content>     & <row\_4\_content>     \\	



		\bottomrule[1.25pt]
	\end{tabular}
\end{table}

Example of Figure Definition

% \begin{figure}[!ht]
% 	\centering
% 	\includegraphics[width=<set_width>]{<file_location>}
% 	\caption[Title of Figure for List of Figures]{Write the figure caption in sentence form. The caption should contain enough detail for the reader to thoroughly understand the figure. So, if the figure is a data plot, be sure to describe the meaning of the various line colors, styles, symbols, etc. used.}
% 	\label{fig:<ref_name>}
% \end{figure}


\chapter{Results \& Discussion}



\chapter{Research Needs}


\chapter{Summary}


\bibliography{UL_general,UL_FSRI,NFPA_std,NIOSH}

\clearpage

\appendix
\captionsetup{list=no}

\end{document}
